% Canonical equation — Fast Fourier Transform (Cooley-Tukey radix-2)
% lege-artis/fourier — v0.1.0
% Companion to fft-cooley-tukey.md
% License: CC-BY-SA-4.0

\documentclass[11pt,a4paper]{article}
\usepackage[utf8]{inputenc}
\usepackage[T1]{fontenc}
\usepackage{amsmath, amssymb, amsthm}
\usepackage[backend=biber,style=numeric]{biblatex}
\addbibresource{../reference-bibliography/refs.bib}

\theoremstyle{plain}
\newtheorem{theorem}{Theorem}
\newtheorem{proposition}{Proposition}
\newtheorem{lemma}{Lemma}

\title{Canonical equation: Cooley-Tukey radix-2 FFT}
\author{lege-artis/fourier}
\date{2026-05-09}

\begin{document}
\maketitle

\section{Setup}

Throughout, $N = 2^m$ for some $m \in \mathbb{N}_0$, and $\omega_N = e^{-2\pi i / N}$
is the primitive $N$-th root of unity.

\section{Decimation-in-time decomposition}

Split $x \in \mathbb{C}^N$ into even- and odd-indexed subsequences:
\[
x_{\mathrm{e}}[r] = x[2r], \qquad x_{\mathrm{o}}[r] = x[2r+1], \qquad r = 0, 1, \ldots, N/2 - 1.
\]

\begin{lemma}[Even/odd split]
\label{lem:split}
With $X[k] = \sum_{n=0}^{N-1} x[n] \omega_N^{nk}$ as in \texttt{dft.tex} Eq.~(1):
\begin{equation}
\label{eq:fft-split}
X[k] = X_{\mathrm{e}}[k] + \omega_N^k \cdot X_{\mathrm{o}}[k], \qquad k = 0, 1, \ldots, N/2 - 1
\end{equation}
where $X_{\mathrm{e}}, X_{\mathrm{o}}$ are length-$N/2$ DFTs of $x_{\mathrm{e}}, x_{\mathrm{o}}$.
\end{lemma}

\begin{proof}
Substitute and use $\omega_N^{2rk} = \omega_{N/2}^{rk}$:
\begin{align*}
X[k] &= \sum_{r=0}^{N/2-1} x[2r] \omega_N^{2rk} + \sum_{r=0}^{N/2-1} x[2r+1] \omega_N^{(2r+1)k} \\
     &= \sum_{r=0}^{N/2-1} x_{\mathrm{e}}[r] \omega_{N/2}^{rk} + \omega_N^k \sum_{r=0}^{N/2-1} x_{\mathrm{o}}[r] \omega_{N/2}^{rk} \\
     &= X_{\mathrm{e}}[k] + \omega_N^k X_{\mathrm{o}}[k]. \qed
\end{align*}
\end{proof}

\section{Periodicity exploit (the second half)}

\begin{lemma}[Second-half identity]
\label{lem:second-half}
For $j \in \{0, 1, \ldots, N/2-1\}$:
\begin{equation}
\label{eq:fft-second-half}
X[N/2 + j] = X_{\mathrm{e}}[j] - \omega_N^j \cdot X_{\mathrm{o}}[j].
\end{equation}
\end{lemma}

\begin{proof}
Setting $k = N/2 + j$ in \eqref{eq:fft-split}, and using
$\omega_{N/2}^{r(N/2)} = (\omega_{N/2}^{N/2})^r = 1^r = 1$ (so
$X_{\mathrm{e}}[N/2+j] = X_{\mathrm{e}}[j]$ and similarly for $X_{\mathrm{o}}$),
together with $\omega_N^{N/2} = e^{-\pi i} = -1$:
\[
X[N/2+j] = X_{\mathrm{e}}[j] + \omega_N^{N/2+j} X_{\mathrm{o}}[j]
         = X_{\mathrm{e}}[j] - \omega_N^j X_{\mathrm{o}}[j]. \qed
\]
\end{proof}

\section{The butterfly}

Combining Lemmas~\ref{lem:split} and~\ref{lem:second-half} yields the
\emph{butterfly operation}: for each $k \in \{0, 1, \ldots, N/2-1\}$,
\begin{equation}
\label{eq:butterfly}
\begin{aligned}
X[k]       &= X_{\mathrm{e}}[k] + \omega_N^k X_{\mathrm{o}}[k] \\
X[k + N/2] &= X_{\mathrm{e}}[k] - \omega_N^k X_{\mathrm{o}}[k].
\end{aligned}
\end{equation}

Each butterfly: one complex multiplication ($\omega_N^k \cdot X_{\mathrm{o}}[k]$)
plus one addition and one subtraction. Two outputs from two inputs sharing one
twiddle factor.

\section{Recursion and complexity}

\begin{theorem}[FFT complexity]
\label{thm:fft-complexity}
With $T(N)$ denoting the cost of computing the DFT of length $N$ via Cooley-Tukey:
\begin{equation}
\label{eq:fft-recurrence}
T(N) = 2 T(N/2) + cN, \qquad T(1) = O(1)
\end{equation}
which by the master theorem yields
\begin{equation}
\label{eq:fft-asymp}
T(N) = \Theta(N \log_2 N).
\end{equation}
\end{theorem}

\begin{proof}
\eqref{eq:fft-recurrence} follows directly from \eqref{eq:butterfly}: solve
two subproblems of size $N/2$, then perform $N/2$ butterflies of constant cost.
\eqref{eq:fft-asymp} is the master-theorem solution for $a = 2$, $b = 2$,
$f(N) = cN$. \cite{Cooley1965}
\end{proof}

Concretely: $\log_2 N$ stages $\times N/2$ butterflies/stage gives
$\frac{N}{2}\log_2 N$ complex multiplications, $N \log_2 N$ complex additions.

\section{Equivalence with direct DFT}

\begin{theorem}[FFT-DFT equivalence in exact arithmetic]
\label{thm:fft-dft-exact}
For $N = 2^m$ and $x \in \mathbb{C}^N$:
\begin{equation}
\mathrm{FFT}(x) = \mathrm{DFT}(x).
\end{equation}
\end{theorem}

\begin{proof}
By induction on $m$. Base case $m = 0$ ($N = 1$): both transforms return $x$
unchanged. Inductive step: assume the equivalence holds at length $N/2$;
Lemma~\ref{lem:split} and Lemma~\ref{lem:second-half} together compute every
$X[k]$ for $k \in \{0, \ldots, N-1\}$ in terms of length-$N/2$ DFTs, which
by hypothesis equal the corresponding length-$N/2$ FFT outputs. \qed
\end{proof}

\begin{proposition}[FFT-DFT equivalence in floating-point]
\label{prop:fft-dft-fp}
In IEEE 754 binary64 arithmetic with rounding-to-nearest:
\begin{equation}
\label{eq:fft-rounding}
\big\| \mathrm{FFT}(x) - \mathrm{DFT}(x) \big\|_\infty
\leq C \cdot \log_2(N) \cdot \varepsilon_{\mathrm{machine}} \cdot \|x\|_\infty
\end{equation}
with $C \leq 4$ a small constant \cite[§12.2]{NumRec3rd}.
\end{proposition}

\eqref{eq:fft-rounding} is the gate condition for the cross-algorithm equivalence
test in \texttt{shared/property-tests/dft.md} (Property~P5). It is a stronger bound
than the direct DFT rounding bound $\sim N \cdot \varepsilon$ because the recursive
halving keeps each addition's roundoff small.

\printbibliography
\end{document}
